\chapter{Projeto de Interface do Usuário}

Na etapa de Engenharia de Software I, a própria equipe produziu a
especificação de requisitos do EduHub, definindo restrições, requisitos
funcionais e não funcionais e histórias de usuário, mas sem chegar a
desenhar \textit{mockups} de interface. Este
capítulo parte desse trabalho anterior: primeiro deriva os princípios de
interface e os fluxos de interação mais relevantes a partir das restrições
e histórias já levantadas e, em seguida, apresenta os protótipos de alta
fidelidade produzidos no Figma para o fluxo mais crítico do sistema, o
caso de uso CU-03.

\section{Princípios de interface}

A interface do EduHub deve seguir os seguintes princípios, todos rastreáveis a
restrições e requisitos não funcionais do documento de requisitos:

\begin{itemize}
    \item \textbf{Aplicação web responsiva} (R-1): o sistema é acessado
    exclusivamente pelo navegador, sem aplicativo móvel nativo nesta versão.
    A interface deve, portanto, se adaptar a diferentes tamanhos de tela,
    especialmente para o perfil Aluno, para o qual a equipe previu maior
    uso em dispositivos móveis ao definir os perfis de usuário na etapa de
    requisitos.
    \item \textbf{Acessibilidade} (R-2 e RNF-10): a interface deve aderir às
    diretrizes WCAG 2.1, nível AA, garantindo o uso do sistema por pessoas com
    deficiência.
    \item \textbf{Intuitividade adaptada ao perfil} (RNF-09): a navegação deve
    ser clara e adequada ao nível de proficiência de cada grupo de usuário.
    Ao definir os perfis de usuário na etapa de requisitos, a equipe
    caracterizou essa necessidade de forma distinta por perfil:
    \begin{itemize}
        \item o \textbf{Aluno} tem baixo nível de acesso, foco em
        visualização de informações e alta expectativa de acessibilidade em
        dispositivos móveis;
        \item o \textbf{Professor} tem nível médio de acesso, foco pedagógico
        e realiza inserção frequente de dados (frequência e notas), exigindo
        uma interface intuitiva que reduza o esforço em tarefas repetitivas;
        \item o \textbf{Administrativo} tem alto nível de acesso e
        conhecimento médio de sistemas de gestão, realizando tarefas de
        cadastro e manutenção de dados mais complexas.
    \end{itemize}
    \item \textbf{Compatibilidade de navegadores} (P-2): a interface deve
    funcionar corretamente em Google Chrome e derivados do Chromium, Mozilla
    Firefox e Safari (versões recentes), sendo esses os únicos navegadores
    contemplados pelo escopo de teste e suporte.
\end{itemize}

\section{Fluxos principais de usuário}

Os fluxos a seguir foram derivados das histórias de usuário e dos casos de
uso da especificação de requisitos. Cada um representa uma tarefa central
para o respectivo perfil, e não uma tela específica.

\subsection{Autenticação e recuperação de senha (HU-01)}

Qualquer usuário do sistema (Aluno, Professor ou Administrativo) precisa
autenticar-se por e-mail e senha antes de acessar suas funcionalidades. Caso
o usuário não recorde a senha, ele deve conseguir solicitar sua recuperação
por meio de um link ou código enviado ao e-mail cadastrado (RF-02). A tarefa
do usuário é, portanto, concluir o acesso ao sistema de forma rápida e
segura. O histórico de login (RF-05) é registrado de forma transparente ao
usuário, sem exigir ação adicional dele.

\subsection{Lançamento de notas e registro de frequência (HU-13 e HU-14)}

Esse é o fluxo central do perfil Professor, que precisa inserir dados com
frequência regular. A partir da lista de turmas e disciplinas vinculadas a
ele (RF-21), o professor deve conseguir lançar e editar notas parciais e
finais por aluno e período (RF-22, RF-23), além de registrar, por aula ou
por data, o status de frequência de cada aluno: presente, ausente ou
justificado (RF-24). O sistema calcula automaticamente o percentual de
frequência (RF-25), sem exigir cálculo manual do professor. Caso o período
letivo já esteja bloqueado pela Coordenação (RF-26), a interface deve
impedir a edição e comunicar claramente esse impedimento, conforme
detalhado no caso de uso CU-03 (Apêndice~\ref{apx:requisitos}). Correções
após o bloqueio dependem de solicitação do professor à Coordenação, que
desbloqueia o período antes de qualquer nova edição.

\subsection{Consulta ao boletim pelo aluno (HU-16)}

O aluno acessa seu boletim para visualizar notas parciais e finais de todas
as disciplinas, organizadas por período letivo (RF-28), e o detalhamento da
sua frequência, incluindo datas de falta e percentual acumulado de presença
(RF-29). Trata-se de um fluxo essencialmente de leitura, alinhado ao perfil
de baixo nível de acesso e foco em visualização definido pela equipe para o
Aluno, e que deve responder rapidamente (RNF-01).

\subsection{Cadastro de usuários e gestão de cargos (HU-03 e HU-04)}

O Administrativo cadastra novos usuários preenchendo os campos obrigatórios
exigidos pelo sistema (RF-13), associando obrigatoriamente um cargo a cada
novo cadastro (RF-06). O mesmo perfil pode alterar dados de terceiros
(RF-11) e inativar usuários mediante confirmação (RF-12). De forma
complementar, o Administrativo pode criar cargos personalizados, definir sua
hierarquia (RF-08) e configurar as permissões associadas a cada cargo
(RF-10), sem poder excluir os cargos padrão do sistema (RF-09). Esse fluxo é
o mais complexo do sistema do ponto de vista de interface, compatível com o
maior nível de acesso e conhecimento médio de sistemas de gestão que a
equipe atribuiu ao perfil Administrativo ao definir os perfis de usuário.

\section{Protótipos de tela}

Para manter o escopo desta iteração enxuto, a prototipação no Figma
concentra-se no fluxo do caso de uso CU-03 (Finalizar e Bloquear Período de
Lançamento), o mais crítico do sistema por proteger a integridade do registro
acadêmico e por girar em torno da regra de trava de lançamento (RF-26). A
Tabela~\ref{tab:telas-figma} lista as telas desse fluxo, do lançamento de notas
e frequência pelo professor até o travamento do período pela coordenação,
associando cada uma ao perfil de usuário correspondente e aos requisitos que a
originam. Cada tela foi prototipada em alta fidelidade em duas versões, uma para
\textit{desktop} e outra para dispositivos móveis, evidenciando a
responsividade exigida pelo princípio R-1. As subseções a seguir apresentam uma
etapa do fluxo por vez, sempre com a versão \textit{desktop} à esquerda e a
versão mobile correspondente à direita, para comparação. Os protótipos dos
demais fluxos e casos de uso serão elaborados de forma incremental nas próximas
iterações.

\begin{table}[ht]
\centering
\caption{Telas do fluxo do CU-03 prototipadas no Figma.}\label{tab:telas-figma}
\begin{tabularx}{\textwidth}{p{0.06\textwidth} p{0.30\textwidth} p{0.16\textwidth} X}
\toprule
\textbf{ID} & \textbf{Tela} & \textbf{Perfil} & \textbf{Fluxo/requisitos relacionados} \\
\midrule
T-01 & Tela de login & Todos & Autenticação por e-mail e senha (HU-01, RF-01) \\
T-02 & Lançamento e edição de notas & Professor & Notas parciais e finais por turma e período, com edição bloqueada quando o período está travado (HU-13, RF-22, RF-23, RF-26) \\
T-03 & Registro de frequência (chamada) & Professor & Presença, falta ou justificativa por aula e cálculo automático de faltas (HU-14, RF-24, RF-25) \\
T-04 & Gestão e travamento de períodos letivos & Coordenação & Verificação de pendências, travamento do lançamento e registro em log de auditoria (CU-03, RF-26, RF-44) \\
\bottomrule
\end{tabularx}
\end{table}

\subsection{O fluxo do CU-03 escolhido}

O CU-03 (Finalizar e Bloquear Período de Lançamento) foi escolhido como fluxo a
prototipar por ser o ponto em que se encontram os dois perfis operacionais do
sistema, o Professor e a Coordenação, em torno da regra que garante a
integridade do registro acadêmico, a trava de lançamento por período (RF-26).
Ele exercita ao mesmo tempo a inserção frequente de dados do Professor (notas e
frequência) e a ação administrativa de encerramento do período, o que permite
demonstrar em um único fluxo tanto o estado editável quanto o estado bloqueado
das telas.

O fluxo, conforme detalhado no CU-03 do Apêndice~\ref{apx:requisitos}, percorre
as seguintes etapas:

\begin{enumerate}
    \item o usuário se autentica na tela de login (T-01), comum a todos os
    perfis;
    \item o Professor acessa suas turmas e, com o período letivo aberto,
    lança e edita as notas (T-02) e registra a frequência da turma (T-03);
    \item a Coordenação acompanha, na gestão de períodos letivos (T-04), o
    progresso de preenchimento de notas e frequência de todas as disciplinas do
    período;
    \item a Coordenação pode bloquear o período a qualquer momento: havendo
    disciplinas com notas ou frequência pendentes, o sistema apenas sinaliza
    essas pendências como aviso, sem impedir o bloqueio, e a Coordenação
    decide se bloqueia mesmo assim, o que altera o status do período para
    ``Bloqueado'' e registra a ação no log de auditoria (RF-44);
    \item a partir desse momento, a trava (RF-26) passa a valer, e o Professor
    que retornar à tela de lançamento encontra os campos desabilitados, sem
    conseguir alterar notas ou frequência do período encerrado;
    \item eventuais correções após o bloqueio dependem de solicitação do
    Professor à Coordenação, que reabre o período pela ação ``Desbloquear
    Período'', realiza os ajustes necessários e bloqueia o período novamente
    ao concluir.
\end{enumerate}

As telas cobrem propositalmente os dois estados do período, aberto e bloqueado,
justamente para tornar visível o efeito da trava, que é o cerne do caso de uso.

\subsection{Autenticação (T-01)}

O acesso ao sistema começa pela tela de login (Figura~\ref{fig:tela-login}),
comum a todos os perfis. O usuário informa e-mail e senha e dispõe do atalho
``Esqueci minha senha'' para a recuperação por e-mail (RF-02). A tela concretiza
a primeira etapa do CU-03 e a história de autenticação HU-01.

\comparativotela{images/tela-cu03-01-login.png}{images/tela-cu03-01-login-mobile.png}%
{T-01 -- Tela de login (desktop à esquerda, mobile à direita), comum a todos os perfis.}{fig:tela-login}

\subsection{Lançamento e edição de notas (T-02)}

Autenticado como Professor, o usuário chega à tela inicial
(Figura~\ref{fig:prof-home}), que lista suas turmas do período corrente, com o
progresso de lançamento de cada uma e atalhos diretos para ``Notas'' e
``Frequência''.

\comparativotela{images/tela-cu03-02-prof-home.png}{images/tela-cu03-02-prof-home-mobile.png}%
{Tela inicial do Professor (desktop à esquerda, mobile à direita), com as turmas do período e atalhos para notas e frequência.}{fig:prof-home}

Na tela de lançamento de notas com o período aberto
(Figura~\ref{fig:notas-aberto}), o Professor edita as notas parciais e finais de
cada aluno por avaliação (RF-22, RF-23). O cabeçalho indica o estado ``Aberto''
do período, e o sistema calcula automaticamente a média a partir da fórmula de
cada aluno, exibida na própria linha. As ações ``Salvar Notas'' e ``Lançar
Notas'' ficam habilitadas.

\comparativotela{images/tela-cu03-03-notas-aberto.png}{images/tela-cu03-03-notas-aberto-mobile.png}%
{T-02 -- Lançamento de notas com o período aberto (desktop à esquerda, mobile à direita): campos e ações habilitados.}{fig:notas-aberto}

\subsection{Registro de frequência (T-03)}

Ainda com o período aberto, o Professor registra a frequência da turma
(Figura~\ref{fig:frequencia-aberto}), marcando cada aluno como presente (P),
ausente (A) ou justificado (J) por aula e data (RF-24). O sistema totaliza as
faltas e calcula automaticamente o percentual de frequência de cada aluno
(RF-25), sinalizando em vermelho quem está abaixo do mínimo, sem exigir cálculo
manual do Professor.

\comparativotela{images/tela-cu03-04-frequencia-aberto.png}{images/tela-cu03-04-frequencia-aberto-mobile.png}%
{T-03 -- Registro de frequência (desktop à esquerda, mobile à direita) com cálculo automático do percentual (RF-25).}{fig:frequencia-aberto}

\subsection{Gestão e travamento de períodos letivos (T-04)}

Autenticada como Coordenação, a usuária chega à tela inicial
(Figura~\ref{fig:coord-home}), que lista os períodos letivos e seu estado
(aberto, com pendências ou fechado), com acesso aos detalhes de cada um.

\comparativotela{images/tela-cu03-05-coord-home.png}{images/tela-cu03-05-coord-home-mobile.png}%
{Tela inicial da Coordenação (desktop à esquerda, mobile à direita), com os períodos letivos e seus estados.}{fig:coord-home}

Ao abrir os detalhes do período (Figura~\ref{fig:periodo-aberto}), a Coordenação
vê, disciplina a disciplina, o progresso de lançamento de notas e frequência,
com o percentual de conclusão de cada uma. Esse painel corresponde à verificação
de pendências prevista no CU-03: o botão ``Bloquear Período'' fica sempre
disponível, e eventuais pendências de notas ou frequência aparecem apenas como
aviso, cabendo à Coordenação decidir se bloqueia o período mesmo assim.

\comparativotela{images/tela-cu03-06-periodo-aberto.png}{images/tela-cu03-06-periodo-aberto-mobile.png}%
{T-04 -- Gestão do período aberto (desktop à esquerda, mobile à direita): progresso por disciplina e ação de bloqueio (RF-26).}{fig:periodo-aberto}

Após o bloqueio, o status do período passa a ``Bloqueado''
(Figura~\ref{fig:periodo-bloqueado}), a ação de bloqueio é registrada no log de
auditoria (RF-44) e a opção disponível passa a ser ``Desbloquear Período'',
executada pela Coordenação quando o Professor solicita correções em notas ou
frequência já encerradas.

\comparativotela{images/tela-cu03-07-periodo-bloqueado.png}{images/tela-cu03-07-periodo-bloqueado-mobile.png}%
{T-04 -- Período após o bloqueio (desktop à esquerda, mobile à direita): status ``Bloqueado'' e opção de desbloqueio.}{fig:periodo-bloqueado}

\subsection{Efeito da trava sobre o Professor}

Fechado o círculo do CU-03, o Professor que retornar à tela de lançamento do
período agora bloqueado (Figura~\ref{fig:notas-bloqueado}) encontra os campos de
nota desabilitados e as ações de salvar e lançar inativas. O cabeçalho informa a
data em que o período foi fechado, comunicando claramente o impedimento previsto
no RF-26. Essa tela é a contraparte da Figura~\ref{fig:notas-aberto} e torna
visível, na interface, o resultado da trava aplicada pela Coordenação.

\comparativotela{images/tela-cu03-08-notas-bloqueado.png}{images/tela-cu03-08-notas-bloqueado-mobile.png}%
{T-02 no estado bloqueado (desktop à esquerda, mobile à direita): campos e ações desabilitados pela trava de lançamento (RF-26).}{fig:notas-bloqueado}

\chapter{Especificação e Estratégia de Testes}

\section{Planejamento e níveis de teste}

A estratégia de testes do EduHub contempla três níveis, cada um com um
objetivo específico dentro do contexto do sistema:

\begin{itemize}
    \item \textbf{Testes de unidade}: verificam isoladamente regras de
    negócio críticas do domínio acadêmico, como o cálculo automático de
    frequência (RF-25), a trava de lançamento por período (RF-26) e a
    prevenção de conflitos de horário (RF-19). Como essas regras envolvem
    lógica de decisão relativamente isolada, o teste de unidade é o nível
    mais eficiente para capturar defeitos nelas o quanto antes.
    \item \textbf{Testes de integração}: verificam a comunicação entre os
    módulos do sistema, por exemplo, entre o módulo de autenticação e o de
    gestão de cargos/permissões, ou entre o módulo de cadastros (turmas,
    disciplinas, professores) e o módulo acadêmico (notas, frequência), que
    depende diretamente dos vínculos criados no cadastro.
    \item \textbf{Testes de sistema}: verificam o comportamento fim a fim das
    histórias de usuário, sob a perspectiva do usuário final, cobrindo os
    critérios de aceitação já descritos em cada história (por exemplo, o
    login em menos de dois segundos, previsto no critério de aceitação da
    HU-01 e no RNF-01).
\end{itemize}

Requisitos considerados críticos para a priorização dos testes são a
autenticação (RF-01, RF-02, RNF-03), o cálculo automático de frequência
(RF-25), a trava de lançamento de notas e frequência (RF-26) e a prevenção
de conflitos de horário (RF-19), por afetarem diretamente a integridade dos
dados acadêmicos e o controle de acesso ao sistema.

\section{Testes de unidade}

As unidades a seguir foram selecionadas por representarem regras de negócio
com lógica de decisão bem delimitada, derivadas diretamente dos requisitos
funcionais:

\begin{table}[ht]
\centering
\caption{Cenários de teste de unidade derivados dos requisitos funcionais.}
\begin{tabularx}{\textwidth}{p{0.16\textwidth} X X}
\toprule
\textbf{Unidade (RF)} & \textbf{Entrada/condição} & \textbf{Resultado esperado} \\
\midrule
Cálculo de percentual de frequência (RF-25) &
Total de aulas registradas e número de faltas de um aluno em uma disciplina &
Percentual de frequência calculado corretamente e exibido sem intervenção manual \\
Validação de conflito de horário (RF-19) &
Tentativa de atribuir um professor ou uma turma a um horário/sala já ocupado por outro professor ou turma &
O sistema rejeita a atribuição e sinaliza o conflito, sem alterar os dados existentes \\
Trava de lançamento por período (RF-26) &
Tentativa de edição de nota ou frequência em um período marcado como ``travado'' &
A edição é bloqueada e uma mensagem informa que o período está encerrado \\
Validação de campos obrigatórios (RF-13) &
Cadastro de novo usuário com um ou mais campos obrigatórios não preenchidos &
O sistema impede o salvamento e indica os campos pendentes \\
\bottomrule
\end{tabularx}
\end{table}

Para unidades que dependem de recursos externos ao próprio módulo, como o
envio de e-mail de recuperação de senha (RF-02) ou de notificações (RF-36),
a equipe pretende utilizar substitutos controlados (\textit{stubs} ou
\textit{mocks}) para isolar a lógica testada do envio real, evitando que a
disponibilidade de serviços externos interfira no resultado do teste.

\section{Testes de integração}

Os módulos do EduHub, autenticação/autorização, cadastros (usuários,
turmas, disciplinas, professores), acadêmico (notas e frequência) e
comunicação (avisos e mensagens), possuem dependências diretas entre si: o
módulo acadêmico depende dos vínculos criados no módulo de cadastros
(RF-17), e todos os módulos dependem da autenticação e das permissões de
cargo (RF-06, RF-10) para autorizar cada ação.

A estratégia adotada é a integração incremental do tipo \textit{bottom-up}:
os módulos de autenticação e de cadastros, por serem a base sobre a qual os
demais módulos operam, são testados primeiro, o que também é coerente com
a ordem que a equipe definiu no planejamento das sprints: a Primeira
Iteração cobre autenticação e cadastro de usuários, e a Segunda Iteração
cobre turmas, disciplinas e alocação de professores. Somente após essa base
estar validada os módulos acadêmico e de comunicação, que dependem dela, são
integrados e testados.

A cada alteração relevante em um módulo já integrado, testes de regressão
são reexecutados sobre os módulos que dependem dele, de forma a identificar
rapidamente efeitos colaterais não previstos. Por exemplo, uma alteração
na regra de permissões de cargo deve reexecutar os testes de integração dos
módulos de cadastros e acadêmico, que dependem dela para autorizar ações.

\section{Critérios de conclusão e cobertura}

Uma funcionalidade é considerada testada de forma satisfatória quando todos
os cenários de teste planejados para os RFs/RNFs associados à respectiva
história de usuário foram executados, o critério de aceitação descrito na
história de usuário correspondente (Apêndice~\ref{apx:requisitos}) foi
verificado, e não há defeito conhecido que bloqueie o critério final da
sprint em que a história está inserida.

A cobertura de código é utilizada como um indicador complementar do esforço
de teste, e não como garantia isolada de qualidade. Um código totalmente
coberto ainda pode conter regras de negócio mal validadas. Como meta
inicial proposta pela equipe, busca-se cobrir prioritariamente as unidades
críticas listadas na Seção 5.2. Essa meta poderá ser revisada ao longo do
projeto.

Além da cobertura, a estratégia de testes se relaciona com as métricas de
produto e de uso que a equipe definiu na etapa de requisitos. Do lado do
produto, o acompanhamento do acoplamento entre classes (CBO), da coesão dos
métodos (LCOM) e da complexidade ciclomática ajuda a identificar trechos de
código mais difíceis de testar e manter. Do lado do uso, a taxa de sucesso,
o tempo de conclusão de tarefas e o questionário SUS (\textit{System
Usability Scale}) complementam os testes automatizados ao avaliar a
experiência real do usuário nos fluxos descritos no Capítulo 4.

\section{Automação e ferramentas}

A execução dos testes de unidade e de integração e a
geração de relatórios de resultado são automatizadas pela esteira de
integração contínua do projeto, o GitHub Actions
(Seção~\ref{sec:build-cicd}), acionada a cada \textit{commit} e
\textit{pull request}, o que garante a reexecução frequente prevista na
estratégia de regressão descrita na Seção 5.3. Para as métricas de produto
definidas pela equipe (CBO, LCOM e complexidade ciclomática), a mesma
esteira executa uma etapa de análise estática de código. As ferramentas específicas
de execução de testes e de análise estática dependem da linguagem e do
\textit{framework} escolhidos para a implementação e são integradas ao GitHub
Actions por meio de scripts de \textit{build} e teste versionados junto ao
código.

\chapter{Gestão de Configuração e Manutenção}

\section{Repositório e itens de configuração}\label{sec:config-repo}

A gestão do Scrumban adotado pela equipe é feita no Taiga, utilizado para
o rastreamento do product backlog e das sprints, e o código-fonte é
versionado no GitHub, que funciona como repositório central do projeto, com o
Git como sistema de controle de versão.

O versionamento no GitHub segue um modelo de branches inspirado no Git Flow,
com dois branches de longa duração: a \texttt{main}, que reflete o código em
produção, e a \texttt{develop}, que reúne o trabalho já integrado e serve de
base para o ambiente de homologação. Cada nova funcionalidade ou correção é
desenvolvida em um branch próprio de curta duração (\texttt{feature/*}),
criado a partir da \texttt{develop}, e retorna a ela por meio de um
\textit{pull request} (PR) revisado. Quando o conteúdo da \texttt{develop}
está estável e validado, ele é promovido para a \texttt{main}, caracterizando
uma nova entrega em produção.

Os seguintes itens de configuração precisam ser controlados ao longo do
projeto:

\begin{itemize}
    \item código-fonte do sistema;
    \item documentação do projeto, incluindo o documento de requisitos (ES1)
    e este documento de projeto (ES2);
    \item modelos e diagramas (diagrama de classes e diagramas de caso de
    uso);
    \item testes (de unidade, integração e sistema, conforme o Capítulo 5);
    \item arquivos de configuração do sistema;
    \item o product backlog e as histórias de usuário, mantidos no Taiga.
\end{itemize}

\section{Gestão de dependências e alterações}

O projeto possui uma dependência externa relevante (D-1): a alocação e
configuração de um servidor de produção pela equipe de TI da instituição de
ensino. Como essa dependência está fora do controle direto da equipe de
desenvolvimento, ela é acompanhada como um risco de cronograma para a fase
de implantação, e não como um item versionado pela equipe.

Para as alterações internas ao sistema, a arquitetura modular exigida pelo
RNF-07, que prevê que novos recursos possam ser adicionados sem impactar
negativamente os módulos existentes, limita o alcance de cada alteração a
um módulo (ou a poucos módulos relacionados), reduzindo o risco de efeitos
colaterais quando uma história de usuário é implementada.

A coerência entre código, testes e documentação é mantida pelo próprio
processo Scrumban adotado pela equipe: o product backlog é um artefato
vivo, refinado continuamente, de modo que qualquer
alteração de escopo em uma história de usuário é refletida no backlog antes
de ser implementada, evitando divergência entre o que está documentado e o
que é efetivamente desenvolvido e testado.

\section{Rastreabilidade}

A rastreabilidade no EduHub segue a cadeia definida na própria estrutura da
especificação de requisitos (Apêndice~\ref{apx:requisitos}): requisitos
funcionais/não funcionais (RFs/RNFs) são agrupados em histórias de usuário,
cada uma listando explicitamente os ``RFs/RNFs associados'', e as histórias
de usuário são alocadas pela equipe às sprints do product backlog. A tabela a seguir
ilustra essa cadeia com dois exemplos concretos:

\begin{table}[ht]
\centering
\caption{Exemplos de rastreabilidade entre requisitos, histórias de usuário e sprints.}
\begin{tabularx}{\textwidth}{p{0.24\textwidth} X p{0.14\textwidth}}
\toprule
\textbf{Requisitos (RFs/RNFs)} & \textbf{História(s) de usuário} & \textbf{Sprint} \\
\midrule
RF-01, RF-02, RF-05 & HU-01 (Autenticação e Segurança) & 1\textsuperscript{a} Iteração \\
RF-15, RF-17 & HU-08 (Criação de Turmas e Disciplinas), HU-09 (Alocação de Professores) & 2\textsuperscript{a} Iteração \\
\bottomrule
\end{tabularx}
\end{table}

\section{\textit{Build}, integração e entrega}\label{sec:build-cicd}

O processo de entrega do EduHub seguirá o modelo Scrumban adotado
pelo projeto: entregas incrementais ao final de cada sprint de 14 dias, cada
uma correspondendo a um incremento funcional do produto, conforme o critério
final que a equipe define para cada sprint no planejamento. Por exemplo, a Primeira
Iteração tem como critério final permitir login e o cadastro de ao menos um
usuário professor. A automação de \textit{build}, integração contínua e
entrega (CI/CD) é realizada no GitHub, por meio do GitHub Actions, apoiada no
modelo de branches descrito na Seção~\ref{sec:config-repo}. A integração
contínua (CI) é acionada a cada \textit{commit} e a cada \textit{pull request},
executando o \textit{build} do sistema, os testes automatizados de unidade e
integração descritos no Capítulo 5 e a análise estática de código. Um
\textit{pull request} só é incorporado à \texttt{develop} quando essa
verificação passa, o que impede que código defeituoso avance. A entrega
contínua (CD) é acionada nas integrações aos branches de ambiente: ao chegar à
\texttt{develop}, a versão é publicada automaticamente no ambiente de
homologação, usado para validação antes da produção, e, ao ser promovida para
a \texttt{main}, é implantada no ambiente de produção.

\chapter*{Controle de Versões do Documento}
\addcontentsline{toc}{chapter}{Controle de Versões do Documento}

Esta seção registra o histórico de alterações significativas deste
documento de projeto. O versionamento do software em si é tratado no
Capítulo 6, na Seção de \textit{Build}, integração e entrega.

\begin{table}[ht]
\centering
\caption{Histórico de versões do documento.}
\begin{tabularx}{0.95\textwidth}{p{0.21\textwidth} p{0.19\textwidth} p{0.12\textwidth} X}
\toprule
\textbf{Data} & \textbf{Autor(es)} & \textbf{Versão} & \textbf{Alteração} \\
\midrule
25/08/2026 & Equipe EduHub & 1.0 & Versão inicial da documentação de projeto (ES2), consolidada a partir do documento de requisitos (ES1). \\
\bottomrule
\end{tabularx}
\end{table}
